\section{Results}
\label{sec:results}

\subsection{Models and Setup}
\label{sec:setup}

We evaluate four models, each with a 5-hour time budget:

\begin{table}[h]
\centering
\small
\caption{Models and agent configurations. Cost is estimated per 5-hour
schedule-only run.}
\label{tab:models}
\begin{tabular}{llccc}
\toprule
\textbf{Model} & \textbf{Agent} & \textbf{Attempts} & \textbf{Cost} \\
\midrule
Claude Code     & Rich CLI (Read/Write/Grep/Bash) & 320 & $\sim$\$20--50 \\
Gemini 2.5 Flash & Rich CLI (Gemini CLI) & 194 & $\sim$\$5--15 \\
Qwen 2.5 Coder 32B & Structured output (vLLM) & 51 & $\sim$\$5 \\
Llama 3.3 70B & Structured output (vLLM) & 8 & $\sim$\$5 \\
\bottomrule
\end{tabular}
\end{table}

Frontier models (Claude, Gemini) use rich CLI agents that read the
codebase before writing schedules. Open-source models (Qwen, Llama)
use a structured-output interface with the skeleton and baseline
schedule injected into the system prompt to provide equivalent
information access.

\subsection{Schedule-Only Results}
\label{sec:results_schedule}

Table~\ref{tab:schedule_results} reports training and held-out AEP for
the best schedule from each model.

\begin{table}[h]
\centering
\small
\caption{Schedule-only results. Baseline: train = 5540.7\,GWh (500
multi-start SGD), ROWP = 4246.7\,GWh. $\Delta$ROWP is the improvement
over baseline on the held-out farm.}
\label{tab:schedule_results}
\begin{tabular}{lcccc}
\toprule
\textbf{Model} & \textbf{Best Train} & \textbf{Best ROWP} & \textbf{$\Delta$ROWP} & \textbf{Key Feature} \\
\midrule
Gemini CLI     & 5552.1 & \textbf{4328.0} & \textbf{+81.3} & Sinusoidal shake \\
Claude Code    & 5600.0 & 4271.5 & +24.8 & Dual Gaussian bumps \\
Qwen 32B       & 5548.8 & TBD    & TBD   & Cosine + coupled $\alpha$ \\
Llama 70B      & 5540.3 & TBD    & TBD   & Exponential decay \\
500-start SGD  & 5540.7 & 4246.7 & ---   & Monotonic decay \\
\bottomrule
\end{tabular}
\end{table}

Gemini's best held-out schedule (+81.3~GWh) was not its best training
schedule. The script with the highest training AEP (5600.0) achieved
only ROWP = 4255.1. This underscores the importance of held-out
evaluation: the training-optimal schedule is not the
generalization-optimal schedule.

\subsection{Discovered Schedule Structures}
\label{sec:discoveries}

The two best-generalizing schedules use qualitatively different
strategies for escaping local optima while maintaining constraint
feasibility.

\paragraph{Claude's dual Gaussian bumps (iter 192).}
The learning rate follows a cosine decay from $4 \times \text{lr}_0$
with two Gaussian bumps centered at $t = 0.5$ and $t = 0.75$:
\begin{equation}
\eta(t) = \eta_{\text{cosine}}(t) + 0.2\,\eta_0\,e^{-\frac{(t-0.5)^2}{2 \cdot 0.04^2}} + 0.3\,\eta_0\,e^{-\frac{(t-0.75)^2}{2 \cdot 0.05^2}}
\end{equation}
The penalty weight dips at $t = 0.6$ between the bumps:
$\alpha(t) = \alpha_{\text{base}}(t) \cdot (1 - 0.5\,e^{-(t-0.6)^2/0.0032})$.
This creates a coordinated explore--enforce cycle: the first bump
escapes a local optimum, the penalty dip permits layout rearrangement,
the second bump explores further, and the final convergence phase
locks in feasibility. Moderate momentum ($\beta_1 = 0.3$,
$\beta_2 = 0.5$) damps oscillations during bumps.

\paragraph{Gemini's decaying sinusoidal shake (iter 67).}
The learning rate is perturbed by a high-frequency sinusoid with
decaying amplitude:
\begin{equation}
\eta(t) = \eta_{\text{cosine}}(t) \cdot \left(1 + 0.08\sqrt{1-t}\;\sin(80\pi t)\right)
\end{equation}
This ``shakes'' the optimizer out of local optima continuously in
early iterations, with the perturbation vanishing at convergence.
The penalty ramp is quadratic and extreme:
$\alpha(t) = \alpha_0 (1 + 7499\,t^2)$, reaching $7500\times$ the
initial penalty at the final iteration. Zero first momentum
($\beta_1 = 0$) eliminates path dependence.

\paragraph{Why these generalize.}
Both schedules share a critical property absent from the baseline:
the penalty weight is \emph{large} at convergence. The baseline
couples $\alpha$ to $1/\eta$, reaching $\sim$100$\times$ at the
end. Claude's schedule reaches $\sim$500$\times$; Gemini's reaches
$7500\times$. On the spacious training polygon (DEI, 50 turbines at
960\,m spacing), even weak constraint enforcement produces feasible
layouts. On the tighter held-out polygon (ROWP, 74 turbines at
792\,m spacing), the extra penalty pressure is essential.

\subsection{Generalization Across Farm Sizes}
\label{sec:generalization}

We evaluate the best schedule from each model on DEI variants with
$N = 30, 40, 50, 60, 70, 80$ turbines.

\begin{table}[h]
\centering
\small
\caption{AEP (GWh) across turbine counts. All LLM schedules are
feasible at every $N$. The seed schedule (single-start, no
multi-start baseline) fails feasibility at most counts (marked *).
Timeout at $N=100$ for full-optimizer scripts. ROWP column uses
the held-out farm ($N=74$, different turbine/polygon).}
\label{tab:generalization}
\begin{tabular}{lccccccc}
\toprule
\textbf{Script} & \textbf{30} & \textbf{50} & \textbf{70} & \textbf{80} & \textbf{ROWP} \\
\midrule
Gemini sched (iter 67)   & 3350 & 5557 & 7728 & 8799 & \textbf{4328} \\
Claude sched (iter 192)  & 3350 & 5560 & 7747 & 8813 & 4261 \\
Qwen full-opt (s1/011)   & 3349 & 5548 & 7709 & 8784 & 4251 \\
Llama full-opt (s2/002)  & 3348 & 5547 & 7707 & 8778 & 4243 \\
Seed (single start)      & 3340 & 5525 & 7698 & 8759 & 4227 \\
\bottomrule
\end{tabular}
\end{table}

The improvement over the single-start seed grows with $N$:
+10~GWh at $N=30$, +32~GWh at $N=50$, +54~GWh at $N=80$.
Schedule-only scripts (Claude, Gemini) consistently outperform
full-optimizer scripts (Qwen, Llama) at every farm size, despite
the full-optimizer scripts having freedom to choose initialization,
multi-start strategy, and hyperparameters.

\subsection{Full Optimizer Mode}
\label{sec:results_full}

In full-optimizer mode, the LLM writes an arbitrary
\texttt{optimize()} function. All four models converge to the same
strategy: wrap \texttt{topfarm\_sgd\_solve} with hexagonal lattice
initialization, standard Adam momentum ($\beta_1 = 0.9$,
$\beta_2 = 0.999$), and higher learning rates (100--200 vs.\
baseline 50). No model writes a custom gradient optimizer that
passes the held-out feasibility check.

\begin{table}[h]
\centering
\small
\caption{Full-optimizer results (best across 5 seeds).}
\label{tab:full_opt}
\begin{tabular}{lcccc}
\toprule
\textbf{Model} & \textbf{Attempts} & \textbf{Best Train} & \textbf{Best ROWP} & \textbf{$\Delta$ROWP} \\
\midrule
Claude Code     & 124 & 5562.5 & 4268.5 & +21.8 \\
Qwen 32B        & 8--11/seed & 5549.5 & 4251.4 & +4.7 \\
Llama 70B       & 3--7/seed & 5546.9 & 4252.0 & +5.3 \\
\bottomrule
\end{tabular}
\end{table}

The full-optimizer mode produces lower held-out gains than
schedule-only for Claude ($+21.8$ vs.\ $+24.8$) and dramatically
lower for Gemini (no comparable run). This suggests that
constraining the action space to the 4-parameter schedule forces
deeper exploration of the optimization dynamics, while the
full-optimizer action space encourages shallow strategy engineering
(initialization tricks, multi-start, hyperparameter sweeps).
